\scanpage{141}

\margin{\deleted{$\unsure{P\mathscr D}$}}

Nous pouvons expliciter le fibré tangent le long des \deleted{fibres} \inserted{la section nulle} de la courbe elliptique universelle rigidifiée à la Torelli, comme étant le fibré formé des $\mathfrak{t}_E$, qui du point de vue fibré vectoriel réel s'identifie au \deleted{fibré} $\mathrm{Sl}(2,\mathbb{Z})$-fibré sur $\mathscr D$ défini par l'opération \unsure{standart|standard} de $\mathrm{Sl},\mathbb{Z})$ sur $\mathbb{R}^2$. Le fibré principal associé \inserted{$P\mathscr D$} s'obtient donc (pour sa seule structure topologique) en remplaçant ds cet énoncé $\mathbb{R}^{2}$ par $(\mathbb{R}^2)^*$. Pour en trouver un revêtement universel, il suffit de choisir un revêtement universel de $(\mathbb{R}^{\unsure{2|1}})^*\simeq\mathbb{C}^*$, p. ex. lé revêtement par l'exponentielle complexe
\[
z\mapsto\exp(2i\pi z):\quad\mathbb{C}\longrightarrow\mathbb{C}^*
\]

% archive sheet number: 140
